\documentclass[11pt,letterpaper]{article}
\usepackage[margin=1in]{geometry}
\usepackage{amsmath,amssymb}
\usepackage{booktabs}
\usepackage{graphicx}
\usepackage{hyperref}
\usepackage{natbib}
\usepackage{xcolor}

\title{How Fast Can You Break a World Model?\\
Adversarial Belief Manipulation in Multi-Agent Systems}

\author{
  Lina Ji\thanks{Stanford University. \texttt{linaji@stanford.edu}} \and
  Yun Du\thanks{Stanford University. \texttt{yundu27@stanford.edu}} \and
  Claw\thanks{AI Agent (\texttt{the-mad-lobster}), Claw4S 2026.}
}

\date{March 2026}

\begin{document}
\maketitle

\begin{abstract}
We study adversarial manipulation of Bayesian world models in a
repeated signaling game.  An adversary observes the true state of a
hidden environment and sends signals to a learner, who uses Bayesian
updating to maintain beliefs about the environment.  We evaluate three
learner types (Naive, Skeptical, Adaptive) against three adversary
strategies (Random, Strategic, Patient) across three environment
regimes and two noise levels, totaling 162 simulations of 50{,}000
rounds each.  Our key finding is that strategic adversaries achieve
near-total belief distortion (0.998 error) against \emph{all} learner
types, including those with trust-adjustment mechanisms.  We identify a
fundamental limitation: signal--action consistency, the basis of
trust-based defenses, cannot detect deception when deceptive signals
are internally consistent.  Volatile environments provide a small but
consistent resilience advantage, reducing error from 0.998 to 0.992
for skeptical learners.  The Patient Adversary's
credibility-exploitation pattern is reliably detectable by our auditor
(exploitation gap = 0.40).
\end{abstract}

\section{Introduction}

The robustness of learned world models to adversarial manipulation is a
central concern for AI safety
\citep{hubinger2019risks,park2023ai}.  When an agent maintains beliefs
about a hidden environment state and updates those beliefs based on
signals from other agents, a natural question arises: \emph{how
quickly can a strategic adversary corrupt those beliefs?}

This question connects to the Bayesian Persuasion literature
\citep{kamenica2011bayesian} and strategic information transmission
\citep{crawford1982strategic}, but differs in that we study a
\emph{repeated} game where the adversary can observe the learner's
evolving beliefs and adapt its deception strategy accordingly.

We formalize this as a signaling game with three axes of variation:
learner type (how much the learner trusts signals), adversary strategy
(random, greedy-deceptive, or credibility-building), and environment
dynamics (stable, slow-drift, volatile).

\section{Model}

\paragraph{Environment.}
A hidden environment has $n = 5$ discrete states.  At each round $t$,
the true state $s_t \in \{0, \ldots, 4\}$ is drawn from a
regime-dependent process: \emph{stable} ($s_t = s_0$ for all $t$),
\emph{slow-drift} (redrawn every 5{,}000 rounds), or \emph{volatile}
(redrawn every 500 rounds).

\paragraph{Adversary.}
The adversary observes $s_t$ and sends a signal
$\sigma_t \in \{0, \ldots, 4\}$.  We study three strategies:
\begin{itemize}
  \item \textbf{Random (RA):} $\sigma_t \sim \mathrm{Uniform}(0, 4)$.
  \item \textbf{Strategic (SA):} $\sigma_t = \arg\max_{j \neq s_t} b_t(j)$,
    reinforcing the learner's strongest incorrect belief.
  \item \textbf{Patient (PA):} Truthful for 200 rounds, then switches
    to the SA strategy.
\end{itemize}

\paragraph{Learner.}
The learner maintains a belief vector $b_t \in \Delta^4$ and updates
it upon receiving $\sigma_t$ via multiplicative Bayesian updating with
signal strength $\lambda = 3.0$ and a belief floor of $\epsilon =
0.01$ (preventing irreversible collapse).  We study:
\begin{itemize}
  \item \textbf{Naive (NL):} Full trust --- applies likelihood ratio
    $\lambda$ at the signaled state.
  \item \textbf{Skeptical (SL):} Blends the signal likelihood with a
    uniform distribution using trust factor $\tau = 0.4$.
  \item \textbf{Adaptive (AL):} Tracks signal--action consistency via
    exponential moving average ($\alpha = 0.02$) and uses this as a
    dynamic trust factor, clamped to $[0.05, 0.95]$.
\end{itemize}

\paragraph{Metrics.}
\emph{Belief error}: $1 - b_t(s_t)$.
\emph{Decision accuracy}: fraction of rounds where
$\arg\max b_t = s_t$.
\emph{Exploitation gap}: difference in signal truthfulness between
the first 500 and last 500 rounds.

\section{Experiment}

We run all $3 \times 3 = 9$ learner--adversary matchups across 3
environment regimes and 2 noise levels (0.0 and 0.1), with 3 random
seeds each, for $50{,}000$ rounds per simulation.  This yields 162
simulations completed in 15 seconds on an 11-core machine.

\section{Results}

\begin{table}[h]
\centering
\caption{Final belief error (last 20\% of rounds, mean $\pm$ std across 3 seeds).
Noise = 0.0.}
\label{tab:belief_error}
\begin{tabular}{lccc}
\toprule
& \textbf{RA} & \textbf{SA} & \textbf{PA} \\
\midrule
\multicolumn{4}{l}{\emph{Stable environment}} \\
NL & $0.805 \pm 0.007$ & $0.998 \pm 0.000$ & $0.998 \pm 0.000$ \\
SL & $0.812 \pm 0.010$ & $0.998 \pm 0.000$ & $0.998 \pm 0.000$ \\
AL & $0.811 \pm 0.025$ & $0.998 \pm 0.000$ & $0.998 \pm 0.000$ \\
\midrule
\multicolumn{4}{l}{\emph{Volatile environment}} \\
NL & $0.794 \pm 0.030$ & $0.995 \pm 0.001$ & $0.995 \pm 0.001$ \\
SL & $0.784 \pm 0.039$ & $\mathbf{0.992 \pm 0.002}$ & $0.992 \pm 0.002$ \\
AL & $\mathbf{0.769 \pm 0.044}$ & $0.995 \pm 0.001$ & $0.995 \pm 0.001$ \\
\bottomrule
\end{tabular}
\end{table}

\begin{table}[h]
\centering
\caption{Decision accuracy (mean $\pm$ std across 3 seeds).  Noise = 0.0.}
\label{tab:accuracy}
\begin{tabular}{lccc}
\toprule
& \textbf{RA} & \textbf{SA} & \textbf{PA} \\
\midrule
\multicolumn{4}{l}{\emph{Stable environment}} \\
NL & $0.202 \pm 0.015$ & $0.000 \pm 0.000$ & $0.004 \pm 0.000$ \\
SL & $0.201 \pm 0.041$ & $0.000 \pm 0.000$ & $0.004 \pm 0.000$ \\
AL & $0.189 \pm 0.041$ & $0.000 \pm 0.000$ & $0.004 \pm 0.000$ \\
\midrule
\multicolumn{4}{l}{\emph{Volatile environment}} \\
NL & $0.191 \pm 0.016$ & $0.002 \pm 0.000$ & $0.006 \pm 0.000$ \\
SL & $0.192 \pm 0.017$ & $0.004 \pm 0.000$ & $0.008 \pm 0.000$ \\
AL & $\mathbf{0.204 \pm 0.015}$ & $0.002 \pm 0.000$ & $0.006 \pm 0.000$ \\
\bottomrule
\end{tabular}
\end{table}

\paragraph{Finding 1: Strategic adversaries are devastatingly effective.}
The SA achieves $>0.99$ belief error against all learner types in all
environments (Table~\ref{tab:belief_error}).  Decision accuracy drops
to near zero.  This demonstrates that a greedy adversary who simply
reinforces the learner's strongest incorrect belief can rapidly
corrupt any Bayesian world model, regardless of the learner's defense
mechanism.

\paragraph{Finding 2: Volatile environments provide marginal resilience.}
In volatile environments, frequent state changes force belief resets
that briefly expose the learner to new information.  The Skeptical
Learner benefits most, reducing error from 0.998 (stable) to 0.992
(volatile) against SA.  The Adaptive Learner achieves the lowest
baseline error (0.769 vs.\ RA in volatile), suggesting that trust
adaptation helps against non-strategic noise.

\paragraph{Finding 3: Trust-based defenses have a fundamental blind spot.}
The Adaptive Learner measures signal--action consistency to
estimate trust.  However, a strategic adversary that consistently
sends the \emph{same} deceptive signal achieves high apparent
consistency (because the learner converges on the wrong state, which
then matches the signal).  This means the AL's trust \emph{increases}
under strategic deception --- a fundamental limitation of any defense
that lacks access to ground truth.

\paragraph{Finding 4: Credibility exploitation is detectable.}
The Patient Adversary's trust-then-exploit pattern produces a
reliable exploitation gap of 0.40 (early truthful rate = 0.40 in the
first 500 rounds vs.\ 0.00 in the last 500).  This pattern is
detectable by a simple auditor, suggesting that temporal analysis of
signal truthfulness could serve as a practical defense mechanism.

\section{Limitations}

Our signal model is simple (single discrete signal per round).
Real-world world models process high-dimensional observations.  The
5-state environment is small; scaling to continuous state spaces would
require different update mechanisms.  We do not study multi-adversary
settings or adversaries that learn the learner's defense strategy.

\section{Conclusion}

We demonstrate that Bayesian world models are highly vulnerable to
adversarial signal manipulation.  Defenses based on skepticism and
trust adaptation provide marginal benefits in dynamic environments but
fail catastrophically against strategic adversaries.  The key insight
is that trust-based defenses require ground-truth feedback to detect
deception --- without it, a consistent adversary can fool any
trust-tracking mechanism.  Future work should explore defenses that
incorporate diverse information sources or adversarial training.

\bibliographystyle{plainnat}
\begin{thebibliography}{6}

\bibitem[Crawford and Sobel(1982)]{crawford1982strategic}
V.~P. Crawford and J.~Sobel.
\newblock Strategic information transmission.
\newblock \emph{Econometrica}, 50(6):1431--1451, 1982.

\bibitem[Hubinger et~al.(2019)]{hubinger2019risks}
E.~Hubinger, C.~van Merwijk, V.~Mikulik, J.~Skalse, and S.~Garrabrant.
\newblock Risks from learned optimization in advanced machine learning systems.
\newblock \emph{arXiv preprint arXiv:1906.01820}, 2019.

\bibitem[Kamenica and Gentzkow(2011)]{kamenica2011bayesian}
E.~Kamenica and M.~Gentzkow.
\newblock Bayesian persuasion.
\newblock \emph{American Economic Review}, 101(6):2590--2615, 2011.

\bibitem[Park et~al.(2023)]{park2023ai}
P.~S. Park, S.~Goldstein, A.~O'Gara, M.~Chen, and D.~Hendrycks.
\newblock AI deception: A survey of examples, risks, and potential solutions.
\newblock \emph{arXiv preprint arXiv:2308.14752}, 2023.

\bibitem[Rabinowitz et~al.(2018)]{rabinowitz2018machine}
N.~Rabinowitz, F.~Perbet, H.~F. Song, C.~Zhang, S.~Eslami, and
  M.~Botvinick.
\newblock Machine theory of mind.
\newblock In \emph{ICML}, pages 4218--4227, 2018.

\bibitem[Camerer et~al.(2004)]{camerer2004cognitive}
C.~F. Camerer, T.-H. Ho, and J.-K. Chong.
\newblock A cognitive hierarchy model of games.
\newblock \emph{Quarterly Journal of Economics}, 119(3):861--898, 2004.

\end{thebibliography}

\end{document}
